
\subsection{Other IPC 2023 Domains}
\label{sec:domains}

We briefly explain the seven domains used in our second experiment with 10
different domains. The other three domains --- Blocksworld, Rovers, and Sokoban
--- where explained in the text. Table~\ref{tab:domain-params-small-exp} shows
the distribution of the main parameters for each domain used in this experiment.

\paragraph{Childsnack}

This domain involves making a number of sandwiches, some with gluten‐free
ingredients for children with allergies, placing them on trays, and delivering
them to tables. Key constraints include production (gluten and non‐gluten
items), packaging (quantity of trays), and serving order.

\paragraph{Ferry}

Cars are distributed across locations, and a ferry (which can fit one car) must
transport them to their target destinations. The domain requires deciding the
ferry’s schedule, routing, and when to move each car, all under capacity
constraints.

\paragraph{Floortile}

A grid of tiles must be painted by robots that can move in four directions via
unpainted tiles. Robots can only move over tiles that are not yet painted, so
they must paint tiles in a certain order.

\paragraph{Miconic}

Elevator domain: there are multiple floors, passengers on origin floors, and
destinations; an elevator must move and pick up / drop off passengers,
respecting boarding and alighting constraints. Though conceptually simple, it
tests planner efficiency regarding sequencing and movement decisions, especially
as passenger numbers and floor counts increase.

\paragraph{Satellite}

Satellites have instruments (which may need calibration) and different modes of
observation; they must orient, turn, switch instruments, and gather images to
satisfy a given goal. Only one instrument can be active at a time, and switching
or calibrating requires actions.

\paragraph{Spanner}

An agent traverses a path where there are $k$ spanners place along this path. At
the end of the path there exists a gate with $m \leq k$ loose nuts. The agent
must collect enough spanners to tighten all nuts and cross corridors to be able
to cross the gate. Spanners break when used, thus the agent must collect enough
spanners to open the gate.

\paragraph{Transport}

This is a traditional logistics problem with capacity constraints. Vehicles must
deliver packages to certain locations on a graph, subject to vehicle capacities
(i.e., number of packages that can be transported at once).

\begin{table}[t!]
\centering
\begin{tabular}{lrl}
\toprule
\textbf{Domain} & \textbf{Param. Dist.} & \textbf{Param. Meaning}               \\ \midrule
Blocksworld     & $n \in [2,10]$        & $n$ blocks to be rearranged           \\
Childsnack      & $c \in [4,12]$        & $c$ children to be served             \\
Ferry           & $a \in [3,30]$        & $a$ cars to be transported            \\
Floortile       & $t \in [2,25]$        & $t$ tiles to paint                    \\
Miconic         & $p \in [1, 10]$       & $p$ passengers to deliver             \\
Rovers          & $r \in [1,4]$         & $r$ rovers to drive and collect items \\
Satellite       & $s \in [1,8]$         & $s$ satellites to synchronize         \\
Sokoban         & $b \in [1,4]$         & $b$ boxes to push to their goals      \\
Spanner         & $k \in [1,20]$        & $k$ spanners to collect               \\
Transport       & $v \in [2,10]$        & $v$ vehicles to drive                 \\
\bottomrule \\
\end{tabular}
\caption{Task distribution and plan lengths for each domain. Showing only the
  main parameter for each domain.}
\label{tab:domain-params-small-exp}
\end{table}

